\scanpage{106}
\mspage{\textcircled{\scriptsize 53}}

\[
\left\{\begin{aligned}
 A(s)&=\coprod_{i\in\Pi}A_i(s)& A_i(s)&=\{t\in S\mid\{s,t\}\in A_i\}\\
 F(s)&=\coprod_{i\in\Pi}F_i(s)& F_i(s)&=\{\{t,t'\}\in\unsure{\mathscr P|\mathcal P|\mathfrak P}_2(S)\mid\{s,t,t'\}\in F_i\}
\end{aligned}\right.\tag*{(14)}
\]

\deleted{et} \unsure{correspondantes|correspondants|correspondante} à une décomposition du contour $P(s)$ en

\[
P(s)=\coprod_{i\in\Pi}P_i(s).\tag*{(15)}
\]

Donc l'axiome de connexité en $s$ ne peut être vérifié que si tous les $A_i(s)$ sauf un seul sont vides ! \deleted{\unsure{Donc|Ceci}} En fait, on ne suppose pas que cet axiome soit satisfait, mais on introduit les

\[
P_i=(S,A_i,F_i,R_{1\,i},R_{2\,i},R_{3\,i}=R_{2\,i}\circ R_{1\,i})\tag*{(16)}
\]

qu'on interprète comme des structures triangulaires sur le \unsure{$\simeq$|$=$} ens. de sommets $S$, moyennant l'axiome ci-dessous

\[
\begin{aligned}
 (\mathrm{T3})\qquad&\forall i\in\Pi,\quad(S,F_i)\text{ est une structure}\\
 &\text{triangulaire i.e. }\bigcup_{f\in F_i}f=S,\ \text{(satisfaisant}\\
 &\text{alors évidemment l'axiome \unsure{corr.|corresp.|conn.} à T2,)}\\
 &\text{et \uline{stricte} en ce sens qu'en plus de}\\
 &\text{(T1) (T2), on a connexité des contours}\\
 &P_i(s).
\end{aligned}\tag*{(17)}
\]

Cela s'exprime donc aussi par les deux conditions
% archive sheet number: 105
